\subsection{EXPONENTIAL OF A LINEAR FORM}

Let V be a finite dimensional vector space, \(\mathbf{S}(V)\) its symmetric algebra. The coalgebra structure of \(\mathbf{S}(V)\) defines on \(\mathbf{S}(V)^{*}\) a commutative and associative algebra structure (Algebra, Chap. III, §11, pp.~579 to 582). The vector space \(\mathbf{S}(V)^{*}\) can be identified canonically with \(\prod_{m\geq0}\mathbf{S}^{m}(V)^{*}\) , and \(\mathbf{S}^{m}(V)^{*}\) can be identified canonically with the space of symmetric m-linear forms on V. The canonical injection of \(V^{*}=\mathbf{S}^{1}(V)^{*}\) into \(\mathbf{S}(V)^{*}\) defines an injective homomorphism from the algebra \(\mathbf{S}(V^{*})\) to the algebra \(\mathbf{S}(V)^{*}\) , whose image is \(\mathbf{S}(V)^{*gr}=\sum_{m\geq0}\mathbf{S}^{m}(V)^{*}\) (Algebra, Chap. III, §11, no. 5, Prop. 8). We identify the algebras \(\mathbf{S}(V^{*})\) and \(\mathbf{S}(V)^{*gr}\) by means of this homomorphism; we also identify \(\mathbf{S}(V^{*})\) with the algebra of polynomial functions on V (Chap. VII, App. I, no. 1).

The elements \((u_m) \in \prod_{m \geq 0} \mathbf{S}^m(\mathrm{V})^*\) such that \(u_0 = 0\) form an ideal J of \(\mathbf{S}(\mathrm{V})^*\) ; we give \(\mathbf{S}(\mathrm{V})^*\) the J-adic topology (Commutative Algebra, Chap. III, §2, no. 5), in which \(\mathbf{S}(\mathrm{V})^*\) is complete and \(\mathbf{S}(\mathrm{V}^*)\) is dense in \(\mathbf{S}(\mathrm{V})^*\) . If \((e_i^*)_{1 \leq i \leq n}\) is a basis of \(\mathrm{V}^*\) , and if \(\mathrm{T}_1, \ldots, \mathrm{T}_n\) are indeterminates, the homomorphism from \(k[\mathrm{T}_1, \ldots, \mathrm{T}_n]\) to \(\mathbf{S}(\mathrm{V}^*)\) that takes \(\mathrm{T}_i\) to \(e_i^*(1 \leq i \leq n)\) is an isomorphism of algebras, and extends to a continuous isomorphism from the algebra \(k[[\mathrm{T}_1, \ldots, \mathrm{T}_n]]\) to the algebra \(\mathbf{S}(\mathrm{V})^*\) .

For all \(\lambda \in \mathrm{V}^*\) , the family \(\lambda^n / n!\) is summable in \(\mathbf{S}(\mathrm{V})^*\) . Its sum is called the exponential of \(\lambda\) and is denoted by \(\exp(\lambda)\) (conforming to Chap. II, §6, no. 1). Let \(x_1, \ldots, x_n \in \mathrm{V}\) ; we have

\[
\langle \exp \lambda , x _ {1} \dots x _ {n} \rangle = \frac {1}{n !} \langle \lambda^ {n}, x _ {1} \dots x _ {n} \rangle = \langle \lambda , x _ {1} \rangle \dots \langle \lambda , x _ {n} \rangle
\]

by Algebra, Chap. III, §11, no. 5, formula (29). It follows immediately that \(\exp(\lambda)\) is the unique homomorphism from the algebra \(\mathbf{S}(\mathrm{V})\) to \(k\) that extends \(\lambda\) .

We have \(\exp (\lambda +\mu) = \exp (\lambda)\exp (\mu)\) for all \(\lambda ,\mu \in \mathrm{V}^{*}\) (Chap. II, §6, no. 1, Remark). Thus, the map \(\exp :\mathrm{V}^{*}\to \mathbf{S}(\mathrm{V})^{*}\) is a homomorphism from the additive group \(\mathrm{V}^*\) to the multiplicative group of invertible elements of \(\mathbf{S}(\mathrm{V})^{*}\) . The family \((\exp \lambda)_{\lambda \in \mathrm{V}^{*}}\) is a free family in the vector space \(\mathbf{S}(\mathrm{V})^{*}\) (Algebra, Chap. V, §7, no. 3, Th. 1).

Lemma 1. Let \(\Pi\) be a subgroup of \(\mathrm{V}^*\) that generates the vector space \(\mathrm{V}^*\) , and \(m\) an integer \(\geq 0\) . Then \(\mathrm{pr}_m(\exp \Pi)\) generates the vector space \(\mathbf{S}^m (\mathrm{V}^*)\) .

By Algebra, Chap. I, §8, no. 2, Prop. 2, any product of \(m\) elements of \(V^{*}\) is a \(k\) -linear combination of elements of the form \(x^{m}\) where \(x \in \Pi\) . But \(x^{m} = m! \operatorname{pr}_{m}(\exp x)\) . Q.E.D.

By transport of structure, every automorphism of V defines automorphisms of the algebras \(\mathbf{S}(\mathrm{V})\) and \(\mathbf{S}(\mathrm{V})^{*}\) ; this gives linear representations of \(\mathbf{GL}(\mathrm{V})\) on \(\mathbf{S}(\mathrm{V})\) and \(\mathbf{S}(\mathrm{V})^{*}\) .

